\section{Week 17}
\sol{17.1}{The constant sequence $x_n=c$ has infinitely many indices inside any positive-radius ball centered at $c$, but only one distinct value. The proof needs infinitely many indices so that it can choose increasing indices. It does not need distinct points; a repeated point is entirely compatible with a convergent subsequence and may make convergence simpler.}
\sol{17.2}{An arbitrary choice could repeat an earlier index or place the selected indices out of order, so the resulting list might not be a subsequence. After choosing $n_{k-1}$, choose $n_k$ from $S_k\setminus\{0,1,\ldots,n_{k-1}\}$. The excluded set is finite while $S_k$ is infinite, so a permissible index remains. This guarantees $n_k>n_{k-1}$ without losing membership in $S_k$.}
\sol{17.3}{For a Cauchy sequence in the discrete metric, apply the Cauchy condition with tolerance $1/2$. Beyond some index, every pair of terms must have distance less than $1/2$, hence distance zero. The sequence is therefore eventually constant and converges in $\NN$, proving completeness. Balls of radius $1/2$ are singletons, so no finite family covers $\NN$; total boundedness fails. The sequence $x_n=n$ has no convergent subsequence because every pair of distinct selected terms remains distance one apart. Completeness controls Cauchy sequences; it does not say every sequence has a Cauchy subsequence.}
\sol{17.4}{Total boundedness asks for a finite cover by balls at \emph{each requested radius}. Fix $r>0$. Since $b-a>0$, choose a positive integer $N$ satisfying $N>(b-a)/r$. This implies the spacing $h=(b-a)/N$ satisfies $0<h<r$.

Use the $N+1$ centers
\[
 c_k=a+kh,\qquad k=0,1,\ldots,N.
\]
The first center is $a$, the last is $b$, and all the others lie between them. If $x=b$, it is already a center. Otherwise, $x\in[a,b)$ lies in one of the finitely many consecutive intervals $[c_k,c_{k+1}]$. Its distance to $c_k$ is at most $h<r$, so it belongs to the radius-$r$ ball about that center. Thus every point is covered.

The list is finite and all its centers belong to the space, as the definition requires. The integer $N$, and therefore the number and locations of centers, may depend on $a,b,r$. Once those are fixed, the same list covers every $x\in[a,b]$; we did not choose a new list after seeing each point.}
\sol{17.5}{The route $v\to u\to w$ makes $w$ two-edge reachable from $v$. However, the additional edge $v\to w$ makes its shortest directed distance one, so $w$ is excluded from the exact second neighborhood. In this graph $N_1^+(v)=\{u,w\}$ and $N_2^+(v)=\varnothing$. Replacing exact second-neighbor membership by mere existence of a two-edge route changes the counted set and can change the truth of an inequality involving its size. It is a semantic change, not a harmless syntax simplification.}
\sol{17.6}{For contraction, a suitable request is: ``Using my knowledge of triangle inequalities and induction, teach the Cauchy condition, completeness, and the geometric-tail estimate needed to prove that repeated applications of this specified contraction reach its unique fixed point. After each concept, show its use in that proof.'' For the subsequence theorem, ask: ``Using the metrics and completeness I already know, explain finite small-ball covers, infinite pigeonhole selection, and increasing subsequence indices, then build the nested-selection argument.'' The first needs a contraction factor and iterative geometric error propagation; the second needs total boundedness and finite-cover extraction. Neither requires the entire collection of concepts in the other.

If you have not run the formal check yourself, a reading record may report the statements and definitions you read, the file and theorem names you inspected, the roles of the dependencies you identified, and any proofs you reconstructed independently. It may also report what the repository's own documentation says about its formal checks, clearly attributed to that source. It may not report that you compiled the project, that its checking procedure passed when you ran it, or that you verified all the imported proofs. Those are separate claims, each needing its own evidence.}
